\appendix
\section*{Appendix F: Category-Theoretic Formalization of the Altitude System}

\subsection*{F.0 Overview}

This appendix presents a full category-theoretic formalization of the
altitude architecture used in the Constellation–Manifold Toolkit.
The purpose is to express:
\begin{itemize}
    \item altitudes as objects of a category,
    \item legal transitions as morphisms,
    \item purity as monoidal identity constraints,
    \item non-generativity as morphism closure under identity,
    \item boundedness as norm-bounded functoriality,
    \item diagnostics as natural transformations,
    \item drift semigroup as an endofunctor,
    \item invariant extraction as a functor into a scalar category,
    \item steering as an automorphism,
    \item collision invariant as a terminal object,
    \item the Series V firewall as an admissible subcategory.
\end{itemize}

\subsection*{F.1 Altitude Category}

Define the altitude category:
\[
\mathcal{C}_{\mathrm{alt}}.
\]

Objects:
\[
\mathrm{Ob}(\mathcal{C}_{\mathrm{alt}}) =
\{A_0, A_1, \Phi_7, D_1, D_2, D_3, \Phi_8, \Phi_9, S_1, I_1\}.
\]

Morphisms:
\[
\mathrm{Hom}(a,b) =
\begin{cases}
\{\tau_{a,b}\}, & T(a,b)=1, \\
\varnothing, & T(a,b)=0.
\end{cases}
\]

Thus the altitude category is a thin category:
\[
|\mathrm{Hom}(a,b)| \le 1.
\]

\subsection*{F.2 Acyclicity as Strict Ranking}

Define a rank functor:
\[
r : \mathcal{C}_{\mathrm{alt}} \to \mathbb{N}
\]
such that:
\[
r(A_0)=0,\; r(A_1)=1,\; r(\Phi_7)=2,\; \dots,\; r(I_1)=9.
\]

For any morphism $\tau_{a,b}$:
\[
r(a) < r(b).
\]

Thus $\mathcal{C}_{\mathrm{alt}}$ is strictly ranked and acyclic.

\subsection*{F.3 Purity as Monoidal Identity Constraint}

Define a monoidal structure:
\[
(\mathcal{C}_{\mathrm{alt}}, \otimes, I)
\]
where $I$ is the identity altitude.

Purity constraint:
\[
P(a) \iff a \otimes I = a.
\]

Thus purity is expressed as:
\[
\forall a,\; \mathrm{id}_a \text{ is the only generative morphism}.
\]

\subsection*{F.4 Non-Generativity as Identity Closure}

Non-generativity requires:
\[
\mathrm{Hom}(a,a) = \{\mathrm{id}_a\}.
\]

Thus no altitude may generate new structure:
\[
\mathrm{Hom}(a,b) = \varnothing \quad \text{unless } T(a,b)=1.
\]

\subsection*{F.5 Boundedness as Norm-Bounded Functoriality}

Let $X$ be the system state space.  
Define a functor:
\[
F : \mathcal{C}_{\mathrm{alt}} \to \mathcal{C}_{X}
\]
where $\mathcal{C}_{X}$ is the category of bounded transformations on $X$.

Boundedness constraint:
\[
B(a) \iff \|F(a)\| < C_a.
\]

Thus:
\[
F(\tau_{a,b}) \text{ is norm-bounded}.
\]

\subsection*{F.6 Diagnostics as Natural Transformations}

Diagnostics $D_1, D_2, D_3$ form natural transformations between functors:
\[
\eta_i : F \Rightarrow F.
\]

For each altitude morphism $\tau_{a,b}$:
\[
F(b) \circ \eta_i(a) = \eta_i(b) \circ F(a).
\]

Thus diagnostics commute with altitude transitions.

\subsection*{F.7 Drift Semigroup as Endofunctor}

Define the drift semigroup:
\[
D : \mathcal{C}_{\mathrm{alt}} \to \mathcal{C}_{\mathrm{alt}}
\]
such that:
\[
D_t \circ D_s = D_{t+s}.
\]

Thus $D$ is an endofunctor satisfying semigroup composition.

\subsection*{F.8 Invariant Extraction as Functor}

Define the invariant extraction functor:
\[
I : \mathcal{C}_{\mathrm{alt}} \to \mathcal{C}_{\mathbb{R}^k}.
\]

For each altitude $a$:
\[
I(a) : X \to \mathbb{R}^k.
\]

Thus invariants are functorial scalar projections.

\subsection*{F.9 Steering as Automorphism}

Define steering:
\[
\Theta : \mathcal{C}_{\mathrm{alt}} \to \mathcal{C}_{\mathrm{alt}}
\]
such that:
\[
\Theta \circ \Theta^{-1} = \mathrm{id}.
\]

Thus steering is an automorphism of the altitude category.

\subsection*{F.10 Collision Invariant as Terminal Object}

Define the terminal object:
\[
I_1 \in \mathrm{Ob}(\mathcal{C}_{\mathrm{alt}}).
\]

For any altitude $a$:
\[
\exists !\; \tau_{a,I_1}.
\]

Thus the collision invariant is the unique terminal altitude.

\subsection*{F.11 Firewall as Admissible Subcategory}

Define the firewall subcategory:
\[
\mathcal{C}_{\mathrm{safe}} \subseteq \mathcal{C}_{\mathrm{alt}}
\]
such that:
\[
\mathrm{Hom}_{\mathrm{safe}}(a,b) =
\begin{cases}
\{\tau_{a,b}\}, & T(a,b)=1 \land P(b) \land N(b) \land B(b), \\
\varnothing, & \text{otherwise}.
\end{cases}
\]

Thus:
\[
\mathcal{C}_{\mathrm{safe}} \text{ is the admissible morphism category}.
\]

\subsection*{F.12 Final Safety Theorem}

\textbf{Theorem.}  
A state $x$ is altitude-safe if and only if its morphisms lie in
$\mathcal{C}_{\mathrm{safe}}$.

Equivalently:
\[
\Omega(x)=1 \iff x \in \mathcal{C}_{\mathrm{safe}}.
\]

This completes the category-theoretic formalization.
